\chapter{UNNECESSARY COMMENTS}

	\section{0001: About Lambdas}
		
		\noindent
		My partner is an old teacher in the university where I studied "informatic engineering". He taught Algebra. In our first meeting, after A LOT OF YEARS after I left my studies, and he left the universitty... my first question for him was: "I need to prove all Irrational Numbers are different". THEN he began to write Irrational infinite decimal symbols as lambdas in the blackboard... and from that day, I began to get the idea to write elements, as chains of lambdas. I put them inside t-uplas, when I began to program it in python: THE POSITION of symbols, needs always to be in the same place. My solution is not arithmetic, is based in coordinates, in positional algorithms, inside abstract infinite data structures, that AFTER, we can transform in well defined functions.
		\\\\
		
		\noindent
		Time after, I understood, thanks to people, that an axiom GUARANTEES you, you can not have the same element inside the same set. Well, better said, you can work with a version of the set WITHOUT repetitions, the axiom says THEY ARE the same set.
		\\\\
		
		\noindent
		My technic works GOOD when elements are DIFFERENT. The more DIFFERENT, the better. I used to say diagonalizations creates my perfect storm, because they guarantee the new element is DIFFERENT :D. But this is too much philosophic. My problems begin when elements are TOO MUCH SIMILAR. :D.
		
	\section{0002: Some facts about CLJAs}
	
		\noindent
		Constructions LJAs are $\pi$-recursive infinite abstract data structures that do not occuppy memory. Better said: infinite memory.
		\\\\
		
		\noindent
		\textless ** \textit{See unnecesary comment 0003, if you want to know WHY I call them $\pi$-recursive} \textgreater
		\\\\
		
		\noindent
		They have the principal property to create a perfect bijection between $\mathbb{N}$ and the LCF language itself defines. It is like a hash function, but instead giving you the position in memory of the chain of LCF, inside the infinite structure... it BUILDS the chain, directly, from a particular natural number. With the inverse flja\_practical function. The direct flja\_practical function, receives a chain, and gives you a natural number. The unique memory, the infinite structure occupies, is the code of boht functions, AND the data structure used to define compositions between CLJAs, through DR paths, and the configurations of each CLJA. For each CLJA used inside the composition, or as a subfunction, we need the space in memory for their functions and constants. And some compositions are "predictable", we don't need to save in memory the functions... just do recursives calls to the same function... with little changes...
		\\\\
		
		\noindent
		CLJAs are the core of all, it is a pity to have to quit them from here, to don't make this document... HUGE. They are the origin of all... but once LCF is created, and the bijections are defined as well defined mathematical functions... the CLJA is not neccesary. They even have the guilt of WHICH design we choose for the language of representations for the original set to study vs $\mathbb{N}$.
		\\\\
		
		\noindent
		Once you have the CLJA... the rest fall by its own weight. It only takes time. It is almost an automatic process. I need to program them, to create tests for them. One CLJA could take me a month... I am getting old... I can teach you all my technics, so you can do it momre faster. Or better. This is why talking about CLJAs takes a lot of time... and it depends in the set we are trying to study. WE NEED TO UNDERSTAND IT VERY WELL. And mathematicians used to define sets that they only know a couple of elements for sure. Without knowing THE SHAPE of the elements, we can not define proper representations.
		\\\\
		
		\noindent
		They are a common pattern between many different famous bijections. Or injectivities. They were created from scratch from a different approach. We can emulate Cantor's Polynomials with them. Or the bijection Gödel used in his famous proof, beteween logical expressions and natural numbers... BUT, we are not going to use potences of prime numbers. The direct and inverse function have the same level of computational complexity and cost. After a year, I was able to convert the direct function in a proper mathematical function, well defined, for ALL CLJAs. The inverse function is still a pure algorithm. I need help with the translation. It took me a year to translate a for that is not a sum or a product. The inverse algorithm has too much ifs, and data structure access, that I don't know how to express properly in mathematics.
		\\\\
		
		\noindent
		They can be used for many different sets, with many different natures, having as cardinal any possible $\aleph$ with any possible natural subindex. AND BEYOND. AND BEYOND THAT BEYOND. Cantor's Theorem is not only false in one particular, very weird case. Like the Russell's Paradox. IT IS FALSE IN ALL ITS INFINITE CASES.
		\\\\
		
		\noindent
		One of their more spectacular properties, is letting us to emulate the cardinal power of the combinations of potence sets. Recursively. With $\mathbb{N}$.
		\\\\
		
		\noindent
		They can have configuration constants. CLJA-FTC, is a composition of two CLJA-PNN, with different configuration in their constants.
		\\\\
		
		\noindent
		Like I am not mathematician, I don't know when people is talking seriously in reddit or not. Some dudes created a set, and they say it has cardinal $\aleph_{\omega}$. The smallest cardinal, bigger than any possible $\aleph$ with natural subindex. To build it, first you need an infinite list of sets, each one having as cardinal, an $\aleph$ with growing natural subindex. After that, you make the union of them all. The cardinal of that infinite union, they said, can not be any possible $aleph_{k}$. Once you try to say, the union has cardinal $\aleph_{k}$, one of the sets inside the infinite union has cardinal $\aleph_{k+1}$. WELL, I CAN BUILD A CLJA FOR A VERSION OF THAT SET... and the potence set of that union... and beyond, and beyond... One of my problems is knowing THE LIMITS of the technic.
		\\\\
		
		\noindent
		I want to test the sentence: "There are sets so big, that they can not be sets". Nothing seems bigger than $\aleph_{0}$ from my point of view. But I ignore a lot of things of set theory.
		\\\\
		
		
	
	\section{0003: $\pi$-recursive}
	
		\noindent
		I called CLJAS $\pi$-recursive as a joke. It is not related to the number. It is related to the sound they make in TV when you say bad words. Why? Because they are FUCKING-RECURSIVE. PEEEEEEEE-RECURSIVE. In spanish, is how the name of the number sounds.
		\\\\
		
		\noindent
		A simple CLJA is an infinite recursive tree-graph. A little more complex than that, in nodes happens things. Nodes are infinite too. Often. They can ESCAPE any level of infinite recursivity, easily, using DR values, without breaking the properties of the previous levels. Those marks that seem useless in LCF. Once you believe you can not be more crazy creating recursivity levels... CLJAs has a property: you can always convert THAT CLJA in a little piece of another bigger CLJA. Like a subfunction of the new CLJA, or making a composition of them. You can change to another level of meaning, using the pop-corn effect, or DR paths. You can deal with infinite chains, using the pop-corn effect, it creates the phenomenom of the system of relations, to being able to deal with infinity chains, using only finite chains. It can work with chains.. and trees. It can deal with infinite unions, and one of the most interesting... Once you think you can not go further... you can emulate the concept "the smallest element, bigger, than all previous ones"... that is used a lot in ordinal numbers. Thanks to it I rediscovered epsilon ordinals without realizing it: I thought I had reached $\omega_{1}$. And the next, and the next, and the next... As I hope you can see in this document... you can emulate potence sets... so we can escalate to any possible aleph. AND BEYOND. Fucking recursive.
		\\\\
		
		\noindent
		I forget it!! A CLJA can be simple node of another CLJA. It is different from being a subfunction of the FLJA of another CLJA ( for the alphabet, once we begin to use chains (LCF), as symbols, to create a new alphabets, an new types of chains). There are CLJAs that are only one node. It is totally crazy. There are some tricks to make them more precise. For example: CLJA-FTC only uses less than 5\% (badly talking) of $\mathbb{N}$. It does not use $LCF_{2c}$, which is the majority of the combinations of LCF. It is a very similar phenomenom than prime numbers inside natural numbers. If I am not wrong you call it density zero in $\mathbb{N}$. Making a risky guess, $LCF_{2c}$ has density ONE inside $\mathbb{N}$. And CLJA-FTC does not use it. The rest of LCF has density ZERO in $\mathbb{N}$. To make more humble $\mathcal{P}(\mathbb{N})$, we only use a little piece of $\mathbb{N}$. We could be more precise, improving the relation, using ALL $\mathbb{N}$, but that is an advanced technic in CLJAs called ``Dinamic Predictible Previous Sets in Multiple Levels''. We eliminate $LCF_{2c}$ from LCF. This idea needs checking.
		\\\\
	
	\section{0004: About the sad history of the name of Complete Relations}
	
		\noindent
		The first thing that many mathematicians do is begin to change names of everything. They don't really contribute anything, just changing names. Before knowing ``this person'' this relations did had no name. I only mention their property to always trying to cover all the domain at any possible cost. But this person, while we were checking MY WORK, called this relations that name, and I liked it. I want to left this clear in case someday he tried to say he colaborated with me... when the unique thing he did, was insult me for finding a mistake in one proof he tried to add, WITHOUT HEARING ME when I said him I had my own solution, and delay the presentation in public of this urgent and important work several years. He watched with his own eyes how I solved him TWO POINTS he called impossible before I solved them in his face. He agreed with the solutions. AND AFTER THAT, he tried to ignore my solution for the third point, and hide our interaction. Dishonest behaviour for a mathematician.
		\\\\
		
		\noindent
		The most funny detail is that his mistake was confusing D-pairs, with normal elements of Domain in one point. Trying to proof TPI is an equivalent of injectivity. TPI solves D-pairs, as a injective function by parts solves subsets of elements. But they are not the same. I knew that much before knowing him. I was stuck in that point for years. The question is HE GOT CONFUSED IN HIS OWN NAMING. But ey!! Who am I? Just the creator!! I din't knew what I was doing... hahahaha... his names were "more clear".
		\\\\
		
		\noindent
		Another guy sended me a paper, where the unique thing he did was changing names, and quitting some properties for TPI ( a stuff we are not going to see here ). And he only create the proofs, aboout his changes, are equivalents of my properties... WITHOUT SOLVING THE REAL PROBLEM: ¿Is TPI a FORMAL RIGOROUS equivalent of injectivity? Quitting some important details for stuff they ignored... they never hear, they never wait until the end... so much prepotence... BUT they did not move the work.... prepotence and cowardy. HE ASKED ANOTHER PERSON to solve a problem I show him... but he never tell anybody about my work or tried to move it.
		\\\\
		
		\noindent
		It was so embarassing.. I gave him the definition of TPI, so two days later he came with an example applied to $\mathbb{Q}$ and $\mathbb{R}$. He said to me: see? I did an example in 48 hours what tooks you 20 years. I needed to say him: it is normal, after I gave you ALL THE DEFINITION AND PROPERTIES DONE. After that he confessed he asked for help for doing it... and it had some little mistakes, that can be fixed easily... in reallity he was proving my point: the technic is general.
		\\\\
		
		\noindent
		The second mistake he commited was saying that he can create a TPI for $\mathbb{Q}$ and $\mathbb{R}$, and like `` we know they don't have the same cardinal, by Cantor's Theorem'', TPI is a wrong tool. We were talking about a TPI applied to $\mathcal{P}(\mathbb{N})$ vs $\mathbb{N}$. HOW THAT CHANGED NOTHING??? Since when the second example has the same cardinal for international community, that he needed to find another one?? AND, I needed to remember him you can not use the conclussion of a theorem to deny a counterexample of it, without other new tool, or judging the counterexample by itself.
		\\\\
		
		\noindent
		They don't deserve NOTHING.
		\\\\
		
		\noindent
		I another hand, there was a dude in reddit, that talked me about infinite interesections. He or she did it to show me DI ( we are not talking about this counterexample, here, neither ) was wrong. WHY PACKs ends empty, because we can not ignore the infinite intersection, after all chops. Years after, that helped me to create TPI, and become crazy some mathematicians :D. He even put me an example, when I was trying to create the concept "group of studies" ( about groups of SNEIs and its maximum $\gamma$ value ). I created two hours in youtube, to solve the problem he found... but I lost his clue. One important piece of those videos, is the system of relations. But the videos are oriented to that old branch. And are in spanish. That is why I need to rewrite this document. In that moment, the concept of relations $r_{\theta_{k}}$ did not existed. In those days I used to say: "A UNIVERSE can solve a group of pairs...". I needed to add the word "relation" and define them. A silly detail, because in the tables, data were separated by coullumns... we only needed to call those coulumns.. relations... instead of Universes. Universes were used BY those relations to solve those D-pairs....
		\\\\
	
	\section{0005: More about Naive CA Th}
	
		\noindent
		NOT EXACTLY A FORMAL PROOF. While all PACKs are disjoint, between them, and not empty, they don't have a singular element in common. We don't have the same element, in two different PACKs. If we apply the axiom of choice, using ALL PACKs, we can choose randomly one element from each PACK, to create a new injective function: all those images are going to be different between them. 
		\\\\
		
		\noindent
		If someone can not do something similar to THIS, by themself, they are trolling. And they are useless to this work. See? I am an intuitive unortodox mathematician. I can not speak with proper rigor, but I try to split it in very simple ideas, so you can easily create formal proofs. I can not do it. It is like: I can not walk to the blackboard, to write the proper symbols... but the colleagues of Steve Hawkings NEVER left him alone, because with a little help, that a lot of people can do, we can enjoy the fruits of the work of that team. See my point?
		\\\\
		
		\noindent
		This proof is not mine, I can not remember who said it in this way, naming the axiom of choice. Annonymous nicks and all that stuffs. The original first two proofs, were made by my partner... They only took him 15 minutes or less. One was just oriented to $\mathcal{P}(\mathbb{N})$ vs $\mathbb{N}$, and the second one to more general cases, for all possible $\aleph$. The third one, that was not a totally formal proof, is in the first paragraph of this section. The fourth one was created by a teacher in ULL... He kept it for him... he was just happy with the result, and I don't understand formal proofs. I am happy when THEY are happy.
		\\\\
		
		\noindent
		In the general cardinal game, oriented to DI, if we find one natural number ( A member of $LCF_{2p}$, translated by the bijection ), in two different image sets (in a D-pair)... what we do is erasing all R-pairs that contain THAT natural number. We eliminate common images, trying to make PACKs disjoint. This is the general idea :D. What we would do in our particular DI, is eliminating ALL its universe too. It's a butchery of R-pairs and elements of $LCF_{2p}$, erradicated completely from the original r, trying to create a $r^{\prime}$ relation where all PACKs are disjoint. It creates a marvellous phenomenom: the smaller are our PACKs, the more disjoint they become... creating the sensation that our original image set is bigger... the smaller it gets ( we erradicate completely its members from ALL the relation)... the bigger it seems... it is so beautifull and weird...
		\\\\
		
		\noindent
		In TPI, thanks to define different relations, we hide the weird concept of ``erradicating R-pairs''. We only show how each universe $\theta_{k}$ is able to solve more and more Families of D-pairs, giving them DISJOINT INFINITE PACKS, throught a concrete relation, well defined.
		\\\\
		
		\noindent
		In DI, our goal is searching all possible combinations of D-pairs from (SNEIs X SNEIs), erradicating all R-pairs that make our PACKs not DISJOINT. After all possible chops, our PACKS must contain, at least, one element inside. Only one empty PACK... destroys our hopes of following the Naive CA th. In TPI we try to reduce the set of Families, not solved with disjoint PACKs (NWSP: Not Well Solved Pairs), to the empty set. The shit-tuation: "S" ends empty or not?. Do you remember it? Both conditions happens EXACTLY, and weirdly, in the same point, under very similar conditions and context. You end discarding all $\mathbb{N}$ (better said: $LCF_{2p}$)... but to do that, you need to discard ALL $\mathcal{P}(\mathbb{N})$ too (better said: $SNEIs$). And that second part is the weird phenomenom that we can interpret in many different ways. We can try to see ITS MEANING in many different ways: TPI; DI; predicting previously all possible combinations of diagonalizations; imitating diagonalizations in an inverse way ( obtaining similar phenomena, changing the roles of the sets ); seeing how the system of relations is really, really ,really, similar to an injective function by infinite parts, but instead solving subsets of elements, it solves subsets of D-pairs (in a growing way); doing one experiment seeing how any possible mathematician ends guaranteeing the exact property that let us say TPI or DI is a practical equivalent of injectivity... What convinced you was TPI: an animation where we can see how subsets from a partition of $\mathbb{N}$, ALONE, are able to solve Families of D-pairs, from an initial case of just two Families, to nothing less than ALL FAMILIES.  
		\\\\
		
		\noindent
		For lost readers, about what is the Shit-tuation:\\
		\textless ** \textit{See unnecesary comment 0006, The SHIT-TUATION} \textgreater
		\\\\
		
		\noindent
		To see the original proof, in a style of ``infinite hotel'':\\
		\textless ** \textit{See unnecesary comment 0007, original proof of Naive CA Th} \textgreater
		\\\\
		
		\noindent
		To see WHY we use PACKs with infinite cardinal:\\
		\textless ** \textit{See unnecesary comment 0008, why we use PACKs with infinite cardinal} \textgreater
		
		\noindent
		These unnecessary comments are like George RR Martin books... you only ask me for a little piece of my work. So much to say. This trick helps me with my TOC. You are loosing so much beautifull and amazing details...
		
	\section{0006: The Shit-tuation}
		
		\noindent
		I am going to transcript, a piece from this thread, from twitter ( you need an account to read more beyond the initial twit of the thread. 27 twits):\\\\
		https://x.com/Fistroman1/status/1956005450911994097\\\\
		In one point of that thread, I begin to explain what are exactly the conditions of the ``Shit-tuation''.
		\\\\
		
		\noindent
		Let's call S to a set with $\aleph_{k}$ cardinality.\\
		"k" could be any possible natural number, even zero.\\
		We are going to chop S many times.\\		
		What is a chop?
		\\\\
		
		\noindent
		Doing a chop in S, means extract from S a concrete list of elements, perfectly defined.\\
		We are going to define chops, by the subset of S that survives into S after the chop, okey?\\
		$S_{i}$ are the elements inside S, not extracted by the chop, $CHOP_{i}$.
		\\\\
		
		\noindent
		Steps and properties we know absolute for sure:\\\\
		1.- We are going to do EXACTTLY $\aleph_{0}$ different chops... no matter the k of $\aleph_{k}$, which will be the cardinality of S.\\\\
		1a) Each chop has a natural subindex. ALL are going to be defined at the same time, with a unique formula.
		\\\\
		
		\noindent
		2.- All chops are independent.\\
		All can be executed in parallel.\\
		All can be executed in an instant.\\
		They can be executed disordered, the order is not important.
		\\\\
		
		\noindent
		3.- Like we are representing each $CHOP_{i}$ witth the subset that survives inside S.\\
		We will have $\aleph_{0}$ "surviving" $S_{i}$ subsets of S, if we study each chop independently.\\
		These are their properties:\\\\
		3a) They are like russian dolls. They are "nested subsets" of S:\\
		$S_{i+1} \subset   S_{i}$\\
		The next one, in index, is strictly contained inside the previous one. This means each chop is better and better extracting members from S. One is able to extract what ALL the previous ones can extract... but a "few" more elements. NOT in quantity... exactly the same entities... and more...\\\\
		3b) The difference between two consecutive, in index, chops, has the same cardinal than S:\\
		$Card (   S_{i} \ S_{i+1}   ) = Card (S)$\\
		None a singular Si has cardinal zero, in the infinite list, but the "improvements", extracting, have the cardinal of S.\\\\
		3c) THE MAGICAL DETAIL:\\
		The infinite intersection of ALL $S_{i}$ is the empty set:\\
		$$ \bigcap_{i=1}^{\infty} S_{i} = \emptyset $$
		\\\\
		
		\noindent
		REMEMBER: We are absolutely sure about all these previous properties. Each one checked by two or more mathematicians. Each $S_{i}$ is smaller and smaller, none is $\emptyset$, but the intersection of them all is $\emptyset$. AND HERE COMES THE WEAK POINT:
		\\\\
		
		\noindent
		\textbf{Question: If we apply ALL chops, at the same time: ¿S ends empty or not?}
		\\\\
		
		\noindent
		Just try to think about it, before reading the answer.
		\\\\
		
		\noindent
		ANSWER:\\
		IT DOES NOT MATTER.\\
		In both cases, Cantor's Theorem is dead.\\
		\\\\
		
		\noindent
		That situation is what I call the "Shit-tuation". Both technics, TPI and DI, use a special set (S: PACKs in DI; NWSP in TPI) to indicate practical equivalence with injectivity. Different in each one, but ending under the same context explained before. With complementary needs. TPI, to be a practical equivalent of injectivity, needs S to be totally empty (in reallity not, but you like stuff easy to decide). DI needs, only, at least, ONE element surviving inside S, to be an equivalent of injectivity.
		\\\\
		
		\noindent
		\textbf{My only weak point, is not knowing the correct answer.}
		\\\\
		
		\noindent
		I don't know if S ends empty or not. AND BELIEVE ME... you could think mathematicians can give an asnwer to this... you need to see how they react once the answer is related to the Cantor's Theorem :D. But in both posibilities, I have a different counterexample that uses it.
		\\\\
		
		\noindent
		To deny TPI, you are FORCED to say S does not ends empty under the conditions of the Shit-tuation. But if you say that, you are saying that in DI, ALL PACKs, do not end empty, after all chops trying to make the relation to meet the conditions of Naive CA Th. If you try to change your mind again, saying PACKs end empty, under the conditions of the Shit-tuation, you are saying NWSP ends empty. So it means relations $r_{\theta_{k}}$, are able to show, how subsets from a partition of $LCF_{2p}$, GROWS until the cardinal of $SNEIs$, measured indirectly, by its hability to give entire Families of pairs (SNEIs X SNEIs), totally disjoint infinite PACKs. From the initial case of two Families, to nothing less than ALL FAMILIES. Because the set of the Families "not well solved", NWSP, ends empty in $\aleph_{0}$ steps, perfectly defined. THERE IS NO ESCAPE. In reallity it's a proof with two paths, where the two paths drive us to the same conclussion: an impossible counterexample of Cantor's Theorem, can be perfectly builded. We just only don't know which one is the correct one :D.
		\\\\
		
		\noindent
		It is so weird. I HAVE NOT FORMAL INJECTIVITIES, because definitions fail us for ignoring these phenomena could exist, BUT I AM ABLE TO MAKE YOU SWEAR, TPI or DI is an injectivity, for the glory of your mother. While you are thinking you are denying TPI or DI. AND YOU, will be an expert. It will be the opinion of an expert. Not mine.
		\\\\
		
		\noindent
		Why I know this would happen? IT HAPPENED, with two groups of different mathematicians across the world, these last 8 years. SO CLEARLY, that I am confident I can repeat it with any possible two groups of any possible mathematicians, under the same roof, being respectfull, and following the rules of one simple experiment, that only takes one morning. Five hours, one morning, three sessions. How to build the system of relations; TPI explained to one group isolated; DI explained to another group isolated. The final step is just make them share notes, and make them realize what they were guaranteeing, while they were thinking they were denying each branch, very sure about themself. They never have patience to hear two branches, and see the consequences of tehir choices, once they believed they have found the mistake.
		\\\\
		
		\noindent
		With Ask Perplexity I don't need all this game. They saw clearly that the decreasing nested subsets of $S_{i}$, their infinite intersection, is not FORCED to be empty. If we can prove it is empty, it has a special meaning. They saw that is not so easy to acomplish the point 3c of the Shit-tuation. So they asked me HOW THE HELL build the relations, to achieve that point. Well, read this document.
		\\\\
			
			
		
		
		

	\section{0007: Original proof of Naive CA Th}
	
		\noindent
		Imagine an infinite cinema. With infinite chairs. We don't know exactly which infinite cardinal have the set of the chairs. It could be any possible $\aleph$.
		\\\\
		
		\noindent
		Suddenly appears an infinite quantity of spectators. We don't know neither, its $\aleph$.
		\\\\
		
		\noindent
		The cinema ushers can not lie. They are magical. Even when they don't know the real answer, or they believe strongly in one wrong answer, their mouth, can not say a lie.	
		\\\\
		
		\noindent
		For COVID rules, in the cinema, each spectator must be seated in the center of a group of three chairs, for their exclusive use. Even leaving one chair distance with the aisle and walls. If S means one spectator, and C means one chair, they must follow this pattern: CSC CSC CSC CSC CSC CSC ... in each line of chairs.
		\\\\
		
		\noindent
		The ushers said to you, THAT EVERYBODY is seated, following perfectly the laws to prevent COVID spread inside cinemas.
		\\\\
		
		\noindent
		QUESTION: ¿Is the cardinal of spectators greater than the cardinal of chairs?
		\\\\
		
		\noindent
		We still IGNORE which $\aleph$ is each one... we only know THAT trustable data.
		\\\\
		
		\noindent
		We don't know if all chairs, or if all lines lines in the cinema, are fully used (respecting the COVID laws)
		\\\\
		
		\noindent
		People normally answer this in seconds: NOPE
		\\\\
		
		\noindent
		This is an imaginary scenary, about how usefull could be a relation, where for each spectator, we are able to add to the relation these three R-pairs:\\
		$(SPECTATOR_{a}, \:\:\: CHAIR_{1})$\\
		$(SPECTATOR_{a}, \:\:\: CHAIR_{2})$\\
		$(SPECTATOR_{a}, \:\:\: CHAIR_{3})$\\
		\\\\
		
		\noindent
		Now read the chapter ``Naive CA Th''. The second paragraph is a formalization of this HISTORY you have just agreed with.
		
	\section{0008: Why we use infinite PACKs}
	
		\noindent
		Irrational numbers, paradox behaviour, emulating them thanks to having infinite representations now, different example of how sets loose properties of order when we create "bijections" wiuth N... they only loose the bijection, but we have another tool very similar. The difference is in properties of order, not in cardinal. 
		\\\\
		
		
		
		

	
	
	
	